\scanpage{292}
On a un isom canonique
\[
T^p(M)=\operatorname{Hom}(T_p(M),\mathbb Z_p)
\]
On notera que a) $T^p$ n'est pas un foncteur exact, mais seulement exact ``mod groupes finis'' b) $T^p$ s'annule sur tous les groupes finis exactement c) De façon précise, $T^p$ établit une \deleted{\unsure{anti-}}\unsure{équiv.} des catégories ``\uline{mod} les groupes finis'', i.e. essentiellement des \unsure{catégories} avec les espaces vecteurs de dim finie sur $\mathbb Q_p$.

Soit $M\longrightarrow N$ un hom de groupes de la forme $\mathbb Q_p/\mathbb Z_p^{\unsure{n}}$, et $H$ un $p$-groupe fini. \uline{Conditions équivalentes}

(i) $T_p(N)\longrightarrow T_p(M)$ injectif à conoyau libre [i.e. isom sur un facteur]

(i bis) $T_p(N)\otimes H\longrightarrow T_p(M)\otimes H$ injectif

(i ter) \deleted{\unsure{$T_p(M)$}}, $T_p(N)/pT_p(N)\longrightarrow T_p(M)/pT_p(M)$ injectif

(ii) $T^p(M)\longrightarrow T^p(N)$ surjectif

(iii) $M\longrightarrow N$ \unsure{épim} sur un quotient facteur direct

(iv) ${}_pM\longrightarrow {}_pN$ est \uline{surjectif}
\mspage{291}
